% Chapter 20: Euler Numbers and Tangent/Secant Expansions
% ch20_euler_numbers.tex

\chapter{Euler Numbers and Tangent/Secant Expansions}

\begin{historical}
Euler discovered in 1755 that the Taylor series of $\tan(x)$ and $\sec(x)$ have
integer coefficients that count alternating permutations. These numbers, now called
Euler zigzag numbers or Andre permutations, appear throughout combinatorics, analysis,
and the theory of continued fractions.

Euler's generating function identity:
\[
\tan x + \sec x = \sum_{n=0}^{\infty} A_n \frac{x^n}{n!}
\]
where $A_n$ is the number of alternating permutations of $\{1, 2, \ldots, n\}$.
This discovery was part of Euler's broader investigation of trigonometric series
and their connections to number theory.
\end{historical}

\section{Prerequisites}

To work through this chapter, you need:
\begin{itemize}
    \item \textbf{Mathematics:} Taylor series, basic combinatorics, understanding of permutations.
    \item \textbf{Programming:} Python functions, iterative algorithms, factorial computation.
\end{itemize}

\section{Euler Zigzag Numbers}

The Euler zigzag numbers $A_n$ count the number of alternating (up-down) permutations
of $\{1, 2, \ldots, n\}$. An alternating permutation satisfies:
\[
a_1 < a_2 > a_3 < a_4 > a_5 < \cdots
\]

Such permutations are also called \textbf{zigzag permutations} or \textbf{Andre permutations}.
The sequence begins:
\begin{align*}
A_0 &= 1 \\
A_1 &= 1 \\
A_2 &= 1 \\
A_3 &= 2 \quad \text{(permutations: 132, 231)} \\
A_4 &= 5 \quad \text{(permutations: 1324, 1423, 2314, 2413, 3412)} \\
A_5 &= 16 \\
A_6 &= 61 \\
A_7 &= 272 \\
A_8 &= 1385 \\
A_9 &= 7936 \\
A_{10} &= 50521
\end{align*}

\subsection{Recurrence Relation}

The zigzag numbers satisfy the recurrence:
\[
2A_{n+1} = \sum_{k=0}^{n} \binom{n}{k} A_k A_{n-k}
\]
This can be proven using the generating function and the fact that $(\tan x + \sec x)' = \sec x (\tan x + \sec x)$.

\begin{lstlisting}[caption={Computing Euler zigzag numbers}, label=lst:euler_zigzag]
from euler import euler_number_zigzag, euler_secant_numbers

print("Euler zigzag numbers A_n:")
for n in range(12):
    print(f"  A_{n} = {euler_number_zigzag(n)}")

print("\nEuler secant numbers E_{2n}:")
print(f"  {euler_secant_numbers(8)}")
\end{lstlisting}

\section{The Generating Function}

Euler's key discovery was the generating function:
\[
\tan x + \sec x = \sum_{n=0}^{\infty} A_n \frac{x^n}{n!}
\]

Separating even and odd powers gives:
\[
\sec x = \sum_{n=0}^{\infty} A_{2n} \frac{x^{2n}}{(2n)!} = 1 + \frac{x^2}{2} + \frac{5x^4}{24} + \frac{61x^6}{720} + \cdots
\]
\[
\tan x = \sum_{n=0}^{\infty} A_{2n+1} \frac{x^{2n+1}}{(2n+1)!} = x + \frac{x^3}{6} + \frac{x^5}{20} + \frac{61x^7}{5040} + \cdots
\]

\subsection{Connection to Bernoulli Numbers}

The coefficients of $\tan x$ can also be expressed in terms of Bernoulli numbers:
\[
\tan x = \sum_{n=1}^{\infty} \frac{(-1)^{n-1} 2^{2n}(2^{2n}-1) B_{2n}}{(2n)!} x^{2n-1}
\]
This shows the deep connection between Euler numbers and Bernoulli numbers, both of
which arise in the study of trigonometric functions.

\subsection{Classical Euler Numbers}

The classical Euler numbers $E_n$ are the even-indexed zigzag numbers:
\[
E_0 = 1, \quad E_2 = 1, \quad E_4 = 5, \quad E_6 = 61, \quad E_8 = 1385, \quad \ldots
\]
By convention, $E_{2n+1} = 0$ for all $n \geq 0$. These are the coefficients in the
Taylor expansion of $\sec x$:
\[
\sec x = \sum_{n=0}^{\infty} \frac{E_{2n}}{(2n)!} x^{2n}
\]

\begin{lstlisting}[caption={Taylor series for tan(x) and sec(x)}, label=lst:taylor_tan_sec]
from euler import tangent_taylor, secant_taylor
import math

x = 0.5
print(f"x = {x}")
print(f"tan({x}) = {math.tan(x):.12f}")
for terms in [3, 5, 7, 10]:
    approx = tangent_taylor(x, terms)
    print(f"  [{terms} terms] = {approx:.12f}  (error: {abs(math.tan(x) - approx):.2e})")

print(f"\nsec({x}) = {1/math.cos(x):.12f}")
for terms in [3, 5, 7, 10]:
    approx = secant_taylor(x, terms)
    print(f"  [{terms} terms] = {approx:.12f}  (error: {abs(1/math.cos(x) - approx):.2e})")
\end{lstlisting}

\section{The Entringer Triangle}

The Entringer triangle provides a recursive way to compute $A_n$.
Entry $E(n,k)$ counts alternating permutations of $\{1, \ldots, n+1\}$
starting with $k+1$ and going down first.

The triangle is constructed as follows:
\begin{itemize}
    \item Row 0: $E(0,0) = 1$
    \item Row $n$: $E(n,k) = E(n,k-1) + E(n-1,n-k)$ for $k = 1, \ldots, n$
\end{itemize}

The leftmost entry of each row gives $A_n = E(n,0)$.

\begin{lstlisting}[caption={The Entringer triangle}, label=lst:entringer]
from euler import entringer_triangle

triangle = entringer_triangle(8)
for i, row in enumerate(triangle):
    print(f"Row {i}: {row}")
\end{lstlisting}

The triangle begins:
\begin{verbatim}
Row 0: [1]
Row 1: [1, 1]
Row 2: [1, 2, 1]
Row 3: [1, 3, 5, 4]
Row 4: [1, 4, 9, 14, 10]
Row 5: [1, 5, 14, 28, 42, 32]
\end{verbatim}

\subsection{Computational Complexity}

The Entringer triangle has $O(n^2)$ entries, and each entry requires $O(1)$ work to compute.
Therefore, computing all zigzag numbers up to $A_n$ takes $O(n^2)$ time and space.
This is efficient for moderate values of $n$ but becomes impractical for very large $n$.

\section{Asymptotic Behavior}

The zigzag numbers grow very rapidly. Euler discovered the asymptotic formula:
\[
A_n \sim 2 \cdot \left(\frac{2}{\pi}\right)^{n+1} \cdot n!
\]
This can be derived from the singularity structure of the generating function
$\tan x + \sec x = \frac{1 + \sin x}{\cos x}$, which has poles at $x = \pm \pi/2, \pm 3\pi/2, \ldots$

\subsection{Verification}

Using Stirling's approximation $n! \sim \sqrt{2\pi n}(n/e)^n$, we can verify the growth rate:

\begin{lstlisting}[caption={Verifying asymptotic growth}, label=lst:asymptotic_growth]
from euler import euler_number_zigzag
import math

print("Verifying $A_n \\sim 2\\cdot(2/\\pi)^{n+1}\\cdot n!$:")
for n in [5, 10, 15, 20]:
    actual = euler_number_zigzag(n)
    asymptotic = 2 * (2/math.pi)**(n+1) * math.factorial(n)
    ratio = actual / asymptotic
    print(f"  n={n:2d}: A_n={actual:>10d}, ratio={ratio:.6f}")
\end{lstlisting}

By $A_{20}$, the value exceeds 37 billion, illustrating the extremely rapid growth.

\subsection{Error Terms in the Asymptotic}

A more precise asymptotic expansion includes correction terms:
\[
A_n = 2 \cdot \left(\frac{2}{\pi}\right)^{n+1} \cdot n! \left(1 - \frac{4}{\pi^2 n} + O\left(\frac{1}{n^2}\right)\right)
\]

\begin{lstlisting}[caption={Refined asymptotic verification}, label=lst:refined_asymptotic]
from euler import euler_number_zigzag
import math

print("Refined asymptotic verification:")
for n in [10, 20, 30, 40, 50]:
    actual = euler_number_zigzag(n)
    leading = 2 * (2/math.pi)**(n+1) * math.factorial(n)
    refined = leading * (1 - 4/(math.pi**2 * n))
    ratio1 = actual / leading
    ratio2 = actual / refined
    print(f"  n={n:2d}: leading ratio={ratio1:.8f}, refined ratio={ratio2:.8f}")
\end{lstlisting}

\section{Combinatorial Interpretations}

\subsection{Down-Up Permutations}

An alternating permutation $a_1, a_2, \ldots, a_n$ satisfies $a_1 < a_2 > a_3 < a_4 > \cdots$
(up-down) or $a_1 > a_2 < a_3 > a_4 < \cdots$ (down-up). Both types are counted by $A_n$.

\subsection{Descent Sets}

The Euler zigzag numbers count permutations with a specific descent set pattern.
A \textbf{descent} in a permutation $\sigma$ is a position $i$ where $\sigma(i) > \sigma(i+1)$.
The number of permutations of $\{1, \ldots, n\}$ with descents at positions $1, 3, 5, \ldots$
is $A_n$.

\subsection{Secant and Tangent Numbers}

The numbers $E_{2n} = A_{2n}$ are called \textbf{secant numbers} (coefficients of $\sec x$),
while $A_{2n+1}$ are called \textbf{tangent numbers} (coefficients of $\tan x$).
Euler showed that these numbers arise naturally in the study of trigonometric integrals.

\section{Euler Numbers and Continued Fractions}

Euler discovered deep connections between Euler numbers and continued fractions. The
continued fraction expansion of $\tan x$ involves only odd powers of $x$:
\[
\tan x = \cfrac{x}{1 - \cfrac{x^2}{3 - \cfrac{x^2}{5 - \cfrac{x^2}{7 - \ddots}}}}
\]
Similarly, $\sec x$ has a continued fraction representation:
\[
\sec x = \cfrac{1}{1 - \cfrac{x^2}{2 - \cfrac{x^2}{2 - \cfrac{x^2}{2 - \ddots}}}}
\]

\subsection{Lambert's Continued Fraction}

Johann Heinrich Lambert (1768) proved that $\tan r$ is irrational for rational $r \neq 0$
using continued fractions. This immediately implies $\pi$ is irrational since $\tan(\pi/4) = 1$.

\begin{lstlisting}[caption={Continued fraction for tan(x)}, label=lst:tan_cf]
from euler import continued_fraction_tan

# Compute convergents of tan(1)
convergents = continued_fraction_tan(1, 10)
print("Convergents of tan(1):")
for i, (p, q) in enumerate(convergents):
    print(f"  [{i}] = {p}/{q} = {p/q:.10f}")
print(f"Actual tan(1) = {math.tan(1):.10f}")
\end{lstlisting}

\section{The Andr\'{e} Permutations Connection}

\subsection{Zigzag to Alternating Bijection}

There is a classical bijection between alternating permutations and certain lattice paths.
The Andr\'{e} permutation of type A counts up-down permutations, while type B counts
symmetric alternating permutations.

\subsection{Secant-Tangent Numbers in Physics}

The Euler numbers appear in quantum field theory as coefficients in the expansion
of Feynman diagrams. Specifically, they count the number of ways to pair vertices
in certain topological configurations.

\section{Extended Combinatorial Applications}

\subsection{Andr\'{e}'s Theorem}

Andr\'{e} (1881) proved that the exponential generating function for alternating permutations
is $\sec x + \tan x$. This result connects several areas:
\begin{itemize}
\item Combinatorics: counting permutations with prescribed ascent/descent patterns
\item Analysis: Taylor series of trigonometric functions
\item Geometry: volumes of certain polytopes related to permutohedra
\end{itemize}

\subsection{Generalized Euler Numbers}

The Euler numbers can be generalized to the \textbf{Andre numbers} $A(n,k)$, which count
alternating permutations of $\{1, \ldots, n\}$ with exactly $k$ rises. These satisfy
a refinement of the basic recurrence and form a triangular array.

\section{Computational Applications}

\subsection{Random Permutation Statistics}

For a random permutation of $n$ elements, the expected number of alternating
subsequences grows as $\frac{4}{\pi^2} n$. This follows from the asymptotic
behavior of Euler numbers.

\subsection{Sorting and Search Algorithms}

The analysis of certain sorting algorithms (like quicksort variants) involves
Euler numbers through the study of permutation patterns. The average number
of comparisons in Shellsort with specific gap sequences relates to zigzag statistics.

\begin{lstlisting}[caption={Analyzing random permutations}, label=lst:rand_perm]
from euler import euler_number_zigzag
import random

def count_alternations(perm):
    count = 0
    for i in range(len(perm)-1):
        if (i % 2 == 0 and perm[i] < perm[i+1]) or \
           (i % 2 == 1 and perm[i] > perm[i+1]):
            count += 1
    return count

# Expected number of alternations in random permutation
n = 10
perms = [random.sample(range(1, n+1), n) for _ in range(10000)]
avg_alts = sum(count_alternations(p) for p in perms) / len(perms)
print(f"Average alternations: {avg_alts:.3f}")
print(f"Expected: {n * 2/math.pi:.3f}")
\end{lstlisting}

\subsection{Euler Numbers in Series Expansions}

Euler numbers appear in many classical series expansions beyond trigonometric functions:

\begin{lstlisting}[caption={Euler numbers in series}, label=lst:euler_series]
from euler import euler_secant_numbers
import math

# Verify sec(x) series coefficients
print("Secant series coefficients:")
for n in range(6):
    E_2n = euler_secant_numbers(n)
    coeff = E_2n / math.factorial(2*n)
    print(f"  x^{2*n}/{(2*n)}! = {E_2n}/{math.factorial(2*n)} = {coeff:.6f}")

# Verify tan(x) series coefficients  
print("\nTangent series coefficients:")
for n in range(5):
    A_2n1 = euler_number_zigzag(2*n+1)
    coeff = A_2n1 / math.factorial(2*n+1)
    print(f"  x^{2*n+1}/{(2*n+1)}! = {A_2n1}/{math.factorial(2*n+1)} = {coeff:.6f}")
\end{lstlisting}

\section{Historical Development}

\subsection{Euler's Original Discoveries}

Euler first investigated these numbers in his 1755 work ``Institutiones calculi differentialis.''
He derived the series for $\tan x$ and $\sec x$ and recognized the combinatorial significance.

\subsection{Andr\'{e}'s Contribution}

D\'{e}sir\'{e} Andr\'{e} (1881) provided the first rigorous combinatorial interpretation,
proving that $A_n$ counts alternating permutations. This established the connection
between analysis and combinatorics.

\subsection{Modern Generalizations}

Modern research has generalized Euler numbers to:
\begin{itemize}
\item \textbf{Secant and tangent numbers of higher order}
\item \textbf{Eulerian numbers} (different concept, counting permutations by descents)
\item \textbf{q-analogues} using combinatorial statistics
\end{itemize}

\section{Modern Relevance}

Euler numbers appear in diverse areas: the coefficients of tangent and secant
functions in Fourier analysis, the enumeration of alternating permutations in
combinatorics, the study of zigzag algebras in representation theory, and in
the analysis of sorting algorithms (specifically, the expected number of
alternating subsequences). The Entringer triangle is also connected to the
Andr\'{e} permutations and the theory of descent sets in symmetric groups.

\section{Frequently Asked Questions}

\begin{itemize}
\item \textbf{Q: What are alternating permutations?} \\
An alternating permutation $a_1, a_2, \ldots, a_n$ satisfies $a_1 < a_2 > a_3 < a_4 > \cdots$
(up-down) or $a_1 > a_2 < a_3 > a_4 < \cdots$ (down-up). Both types are counted by $A_n$.

\item \textbf{Q: Why are odd-indexed classical Euler numbers zero?} \\
The classical Euler numbers $E_n$ are defined as the coefficients of $\sec x$. Since $\sec x$ is an even function, only even powers appear. Thus $E_{2n+1} = 0$.

\item \textbf{Q: How fast do Euler numbers grow?} \\
$A_n$ grows very rapidly: $A_n \sim 2 \cdot (2/\pi)^{n+1} \cdot n!$. By $A_{20}$, the value exceeds 37 billion.

\item \textbf{Q: What is the connection to Bernoulli numbers?} \\
Both sequences appear in trigonometric Taylor series. $\tan x$ coefficients relate to Bernoulli numbers, while $\sec x$ coefficients are the Euler numbers. They are both "zigzag" in nature.

\item \textbf{Q: Can the Entringer triangle be used for computation?} \\
Yes. Each row is built by cumulative sums of the previous row, alternating direction. This is known as the Seidel construction and allows efficient computation of all $A_n$ up to a given bound.

\item \textbf{Q: What is the relationship between Euler numbers and the secant/tangent functions?} \\
The Taylor series coefficients of $\sec x$ are the Euler (secant) numbers $E_{2n}/(2n)!$, while the coefficients of $\tan x$ are the tangent numbers $A_{2n+1}/(2n+1)!$.

\item \textbf{Q: How are Euler numbers used in practical applications?} \\
Euler numbers appear in the analysis of algorithms (expected number of alternations in random permutations), in quantum field theory (Feynman diagram combinatorics), and in the study of special functions.

\item \textbf{Q: What is the relation to combinatorial geometry?} \\
The Euler numbers count the number of regions in certain hyperplane arrangements and relate to volumes of permutahedra and associahedra.

\item \textbf{Q: What is the difference between Euler zigzag numbers and classical Euler numbers?} \\
The zigzag numbers $A_n$ include both even and odd indices. The classical Euler numbers $E_n$ are defined as $A_{2n}$ with $E_{2n+1} = 0$.
\end{itemize}

\section{Exercises}

\begin{exercise}
Prove that $A_n$ satisfies the recurrence $2A_{n+1} = \sum_{k=0}^{n} \binom{n}{k} A_k A_{n-k}$ for $n \geq 1$.
\end{exercise}

\begin{exercise}
Compute the first 12 Euler zigzag numbers and verify them by enumerating all alternating permutations of $\{1, \ldots, 6\}$.
\end{exercise}

\begin{exercise}
Use the Taylor series for $\tan x$ to approximate $\tan(0.1)$ to 10 decimal places. How many terms are needed?
\end{exercise}

\begin{exercise}
Show that the Entringer triangle satisfies the recurrence $E(n,k) = E(n,k-1) + E(n-1,n-k)$ with $E(n,0) = E(n-1,0)$.
\end{exercise}

\begin{exercise}
Prove that $A_n = \sum_{k=0}^{n-1} \binom{n-1}{k} A_k A_{n-1-k}$ for $n \geq 2$ by considering the position of the largest element in an alternating permutation.
\end{exercise}

\begin{exercise}
Using the generating function $\tan x + \sec x$, derive the first 8 terms of the Taylor series for $\sec x$ and $\tan x$ by hand.
\end{exercise}

\begin{exercise}
Show that the number of down-up permutations of $\{1, \ldots, n\}$ equals the number of up-down permutations.
(Hint: Consider the complement permutation $\sigma'(i) = n+1 - \sigma(i)$.)
\end{exercise}

\begin{exercise}
Verify numerically that $A_n \sim 2 \cdot (2/\pi)^{n+1} \cdot n!$ for $n = 10, 15, 20$.
Compute the ratio $A_n / [2 \cdot (2/\pi)^{n+1} \cdot n!]$ and show it approaches 1.
\end{exercise}

\begin{exercise}
The secant numbers $E_{2n}$ satisfy the recurrence $E_{2n} = \sum_{k=0}^{n-1} \binom{2n}{2k} E_{2k} E_{2n-2k-2}$ for $n \geq 1$ with $E_0 = 1$. Prove this recurrence.
\end{exercise}

\begin{exercise}
Show that $\int_0^{\pi/2} \sec^2 x \, dx$ diverges while $\int_0^{\pi/4} \sec^2 x \, dx = \sqrt{2} - 1$. Use this to explain why the Taylor series for $\tan x$ has radius of convergence $\pi/2$.
\end{exercise}

\begin{exercise}
Using the generating function, prove that $\sum_{n=0}^{\infty} A_n \frac{x^n}{n!} = \tan x + \sec x$ by showing that $y' = 1 + y^2$ where $y = \tan x + \sec x$.
\end{exercise}

\begin{exercise}[Computational]
Implement the Entringer triangle construction and use it to compute $A_n$ for $n$ up to 20. Compare results with the recursive formula.
\end{exercise}

\begin{exercise}
Prove that the continued fraction coefficients of $\tan x$ are related to Euler numbers by the recurrence relation derived from the generating function.
\end{exercise}

\begin{exercise}
Show that $\sec x + \tan x = \tan\left(\frac{\pi}{4} + \frac{x}{2}\right) \sec x$. Use this to derive alternative series expansions.
\end{exercise}

\begin{exercise}
Using the relation between Euler numbers and Bernoulli numbers, express $A_{2n+1}$ in terms of $B_{2n}$ and verify for small $n$.
\end{exercise}
